% r03_hyb_walk: proof of \Cref{rot:prop:BD}, part V: the random walk with periodic cells, the boundary clause, the remaining kernels, the slice clauses, the levels 0 and 1; this completes part (a).

\begin{proof}[Proof of \Cref{rot:prop:BD}(a)]
\emph{Random walk.} Let $\Omega$ be $\Lambda_X$ or a torus fitted by both lattices with $P\subset M\Z^3$, $A$ regular on $\Omega$. For $j\in\Z^3$ let $\square_j$ be the periodic cell of side $2M$ ``centred'' at $Mj$ and fitted by both lattices (which is possible iff $25\mid2M$, as the tori of the chain allow; \cite[p.~345]{K2}: ``blocks $\{\square_j\}$ which are small scale versions of $D_J$, that is they can be fitted by both the usual lattice and the rotated lattice''), $h_j$ the partition of unity of \cite[\S2]{Ba4} ($h_j=1$ on the $M$-cube around $Mj$, supported in the $\frac43M$-cube, $\sum_jh_j^2=1$), and $\tilde A_j=A_0^{(j)}+\theta_jA'^{(j)}$ with $A_0^{(j)}=A(Mj)$ and the cutoff $\theta_j$ of \cite[p.~575]{Ba4}, so that $\tilde A_j=A$ on $\supp h_j$, is regular with constant $c'=O(Mc)$, and is constant near the boundary of the fundamental region of $\square_j$. With $G_0\defeq\sum_jh_jG^X_r(\square_j,\tilde A_j)h_j$ one computes, exactly as \cite[(2.3)--(2.9)]{Ba4} and \cite[(5.8)--(5.12)]{K2},
\[
(-\Delta_A+m_r^2+a_rP_X(A))G_0=\mathbf 1-\cR,\qquad \cR=\sum_jK_jG^X_r(\square_j,\tilde A_j)h_j,
\]
\[
(K_j\phi)(x)=\mathbf 1_{\square_j}(x)\Big[\sum_{b\in\mathrm{st}(x)}(\partial h_j)(b)(D_A\phi)(b)+(\Delta h_j)(x)\phi(x)-(R_r(A,\partial h_j)\phi)(x)\Big],
\]
where $R_r(A,\partial h)\phi(x)=\sum_{x'\in B(z(x))}\eta'^3U(A(\Gamma_{z,x}))U(A(\Gamma_{z,x'}))^{-1}(\partial h)(\Gamma_{x,x'})\phi(x')$ collects the commutator of the pin with the multiplication by $h_j$ (\cite[(2.8)]{Ba4}); $(\partial h)(\Gamma_{x,x'})=h(x')-h(x)$ for any path, so $|R_r(A,\partial h_j)\phi(x)|\le\sup|\partial h_j|\cdot2.1\cdot\sup_B|\phi|$ by \Cref{rot:hyb:lem:len}(i). The three identities used --- the product rule for $-\Delta_A$ (\cite[(2.5)]{Ba4}), the commutator of the pin with a smooth function (\cite[(2.8)]{Ba4}), and the agreement of the cell operator with the global one on $\supp h_j$ (\cite[(2.6)]{Ba4}; for a periodic cell this holds because $\supp h_j$ stays at distance $\ge\sqrt3$ from the boundary of the fundamental region, so no block and no bond of $\supp h_j$ wraps) --- hold for any block system with contours from the root to the point. Hence $G^X_r(\Omega,A)=\sum_{n\ge0}G_0\cR^n$ in $L^2$ (\Cref{rot:hyb:cor:L2} gives $\|\cR\|_{2\to2}\le O(M^{-1})$ as in \cite{Ba4}, since $\sup|\partial h_j|=O(M^{-1})$ and $\sup|\Delta h_j|=O(M^{-2})$), i.e.\ the path representation \cite[(2.13)]{Ba4}, \cite[(5.56)]{K2}.

\emph{Bounds.} \cite[(2.18)--(2.22)]{Ba4} verbatim: for $x,x'$ in $\supp h_{\omega_0}$ and $f$ supported in $\square_{\omega_n}$, each factor $\|K_{\omega_i}G^X_r(\square_{\omega_i},\tilde A_{\omega_i})h_{\omega_i}\|_{\infty\to\infty}$ is $\le O(M^{-1})c_1$ by \Cref{rot:hyb:lem:HOL} and the bound on $K_j$ ($c_1$ the uniform $\|\cdot\|_{1,0}$ constant of \Cref{rot:hyb:lem:HOL}: $K_j$ takes one covariant derivative of its argument, so the choice of $M$ below is uniform in $r$), the first factor is bounded in $\|\cdot\|_{1,\alpha}$ by \Cref{rot:hyb:lem:HOL} (the only place where $C_\alpha\log^{3\alpha}(1/\eta')$ enters, and only in the clause (3.8)), the number of paths of length $n$ is $\le2^d(3^d)^{n-1}2^d$, and $n\ge M^{-1}\dist(\{x,x'\},\supp f)-2$; for $M$ fixed with $3^dc_1O(1)M^{-1}\le e^{-1}$ the sum converges and gives \cite[(3.7)--(3.8)]{K1} with $\delta_0=(2M)^{-1}$. The tail with $n>n_0$ is handled with the $L^p$ bounds (2.17) as in \cite[(2.21)]{Ba4}. The bounds (3.6) for $\Delta^{(r)}$ and $C^{(r)}$ are \Cref{rot:hyb:prop:CT} and \Cref{rot:hyb:cor:Ck}; the bounds for the next-step precision $\Delta^{(r)}+aL^{-2}\hat Q^*\hat Q$ and its restricted inverses $\hat C^{(r)}_\Lambda$ are \Cref{rot:hyb:prop:CLam}.

\emph{Two-torus clause.} For two domains $\Omega_1$, $\Omega_2$ sharing a fundamental region $F$ as in \Cref{rot:hyb:lem:domdiff} (or a torus against $\Lambda_X$), the representations coincide term by term for the paths whose cells do not meet the boundary of $F$, and the difference is bounded, as on \cite[p.~579]{Ba4}, by the same sum restricted to paths meeting the boundary; such a path of $n$ cells also joins $x$ to $\supp f$ in its own domain, so $\sqrt3Mn\ge\max\{d_i(\{x,x'\},\supp f),\dist(\{x,x'\},\partial F)+\dist(\supp f,\partial F)\}-c$ for $i=1,2$, and the average of the two exponents gives the factor $e^{-(2\sqrt3M)^{-1}(d_\wedge(\{x,x'\},\supp f)+\dist(\{x,x'\},\partial F)+\dist(\supp f,\partial F))}$ for the difference of the $G^X_r$ clauses (3.7)--(3.8), with $d_\wedge\le d_i$ the smaller of the two quotient distances. For $\Delta^{(r)}$, $C^{(r)}$ and $\hat C^{(r)}_\Lambda$ the two-torus clause is \Cref{rot:hyb:lem:domdiff} and \Cref{rot:hyb:prop:CLam}. The metric $d_\wedge$ cannot be replaced by the distance of the lift (the remark after \Cref{rot:hyb:lem:domdiff}); the uses of the clause in this section (\Cref{rot:lem:CB}(N), the kernel sweep of \Cref{rot:thm:Asrc}, and \Cref{rot:lem:TC}(IV)) need only the decay on each torus in its own quotient distance together with the boundary factor, which is what is proved.

\emph{The remaining kernels.} $H^X(A)=a_rG^X_r(A)Q_X(A)^*$: apply (3.7)--(3.8) to $f=w_z$ ($\|w_z\|_\infty=1$, support in a block of diameter $\le\sqrt3$). $B^X(A)=a_r^{-1}+Q_X(A)C_AQ_X(A)^*$: by the diamagnetic inequality $|C_A(x,y)|\le C_0(x,y)=(-\Delta+m_r^2)^{-1}(x,y)$ entrywise ($-\Delta_A+m_r^2=c-T_A$ with $|T_A|=T_0$ entrywise and $(c-T_0)^{-1}=\sum c^{-n-1}T_0^n\ge0$), so $|B^X(z,w)|\le a_r^{-1}\delta_{zw}+\sum_{x\in B(z),y\in B(w)}\eta'^6C_0(x-y)\le C(1+|z-w|)^{-1}$ ($d=3$). $\Gamma^X(A)=H^XB^XH^{X*}$: the fine-variable regularity is that of $H^X$, and $|\Gamma^X(x,y)|\le\sum_{z,w}Ce^{-\delta_0|x-z|}(1+|z-w|)^{-1}e^{-\delta_0|y-w|}\le C(1+|x-y|)^{-1}$. $B^X(A)$ and $\Gamma^X(A)$ are not unit-decaying: $\Gamma^X(x,y)\approx C_A(x,y)$ for $|x-y|\gg1$, at the coarse mass $m_r=m\delta_r\to0$; \cite[Thm.~3.3]{K1} never claims otherwise (it concerns $G_r$, $\Delta^{(r)}$, $C^{(r)}$). Where a renormalised family carries a line $\Gamma^F$ or $(A^F)^{-1}$ at the top scale, either the coarse-slice decomposition of \Cref{rot:cor:unif}(b) is used, or, in the localised form \cite[(3.55)]{K1}, all vertex positions lie in the region $F_v$ of \Cref{rot:lem:CB}, of diameter $\le Cp_r$, where a line bounded by $C(1+|x-y|)^{-1}$ costs at most a factor $Cp_r^3$, absorbed in the field-norm polynomial of \Cref{rot:prop:cp}.

\emph{The slice clauses.} Nothing in the argument above uses the level: the same proof at level $j\le r$, in the unit scale of level $j$, gives (3.6)--(3.8) for the step-$j$ operators $G_j$, $\Delta^{(j)}$, $C^{(j)}$, $H_j$ of the hybrid flow with constants uniform in $j$ (the first step, with the tiles as blocks, is covered by \Cref{rot:hyb:lem:coer} and \Cref{rot:hyb:lem:A0} with $B^X_1$ in place of $B^X_r$; the steps $j\ge2$ are the standard rotated steps acting on composite blocks, \Cref{rot:hyb:lem:alg}(iv)). King's derivation of the slice bounds \cite[(3.63)--(3.65)]{K1} from Theorem~3.3 (\cite[Prop.~3.7]{K1}: the slice $S_j=H_jC_jH_j^*$ at a regular background, the one- and two-ended contour clauses summed along contours of bounded length) therefore applies verbatim to the hybrid flow, with the covariant lattice derivatives and transported H\"older quotients of \cite[(2.10)--(2.11)]{Ba3} at regular $A$; after the constant-background reduction of \Cref{rot:lem:CB} the compared lines are the $A=0$ slices of \Cref{rot:hyb:lem:A0}, so there the plain form holds. For the contour clauses the contour lengths are those of \Cref{rot:hyb:lem:len}(i).

\emph{Levels $0$ and $1$.} In the form in which the hybrid lower bound (\Cref{rot:prop:H4}) consumes (a) at its first steps (precision bounds, precision decay, sup-decay of $C^{(0)}$ and $C^{(1)}$), (a) holds also at $j\in\{0,1\}$ by the same proofs, which use $r\ge2$ nowhere: at $j=1$ the level-$1$ block is the tile, inside the ball of radius $\frac{\sqrt3}2$ about its centre, with an in-tile contour of length $\le1$ in the unit scale of level $1$, so \Cref{rot:hyb:lem:coer}, \Cref{rot:hyb:cor:L2}, \Cref{rot:hyb:prop:CT} and \Cref{rot:hyb:lem:domdiff} read verbatim (the first Neumann eigenvalue of the cube is then $100\sin^2(\pi/40)>\frac14$), and the fine spacing $\eta'=\frac15$ is a constant; at $j=0$ the precision $-\Delta_A+m_0^2+aL^{-2}\hat Q^*\hat Q$ has finite range and its lower bound is the $(r+1)$-pin clause of \Cref{rot:hyb:cor:L2} at $r=0$.

\emph{The Neumann variant.} \cite[Thm.~3.3]{K1} is printed for Neumann boundary conditions on a rectangular parallelepiped of large blocks. For the hybrid flow such a domain is a digitised rotated parallelepiped, on which neither the reflection representation \cite[(2.42)]{Ba4} nor the axis-cube cover near the boundary is available; the Neumann variant is not claimed and not used: \cite[\S\S4--5]{K2} decouple with periodic cells in the random walk, the localisation of \Cref{rot:lem:CB} runs on the infinite lattice in the universal-cover form, and the only Neumann object of the chain, the boundary factor of $\delta G_r$, is replaced by the two-torus clause above and the same-domain clause \Cref{rot:lem:CB}(N).

\emph{The standard flow.} For the flow $Y$ every statement of (a) is \cite[Thm.~3.3]{K1}, proved in \cite{Ba4} on $\eta'\Z^3$ in the frame of $\Yk$, with King's uniform constants; its two-torus clause and restricted inverses are \cite[Prop.~5.6]{K2} and \Cref{rot:hyb:lem:domdiff}, whose proof uses nothing specific to the hybrid blocks. This completes the proof of (a).
\end{proof}
